% =============================================================================
% L08 (Class 8) --- textbook exercises with solutions. Built from the MPE list template
% (~/.claude/commands/templates/lista-exam.tex).
%
% TWO MODES (one source file):
%   \solucoesfalse  -> blank LIST (space to solve)
%   \solucoestrue   -> ANSWER KEY (same file)
%
% Compile (copyable PDF, global lmodern+cmap rule):
%   pdflatex lista-L08-book-exercises.tex  (run 2x) ; check: pdffonts ... (all Type 1)
% Numbers: python Listas/L08-book-exercises-check.py reproduces and asserts every number.
% =============================================================================
\documentclass[11pt,a4paper]{article}
\usepackage[utf8]{inputenc}
\usepackage[T1]{fontenc}
\usepackage{lmodern}   % Type 1 vector fonts (replaces PK bitmaps)
\usepackage{cmap}      % ToUnicode CMap -> copyable text and accents
\usepackage[english]{babel}
\usepackage{amsmath,amssymb,amsfonts}
\usepackage[a4paper,margin=2.2cm]{geometry}
\usepackage{enumitem}
\usepackage{fancyhdr}
\usepackage{comment}
% One cut file per job: parallel builds in one folder would otherwise share comment.cut and
% swap solutions. Redefine DefaultCutFileName; overriding CommentCutFile alone is reset by comment.sty.
\def\DefaultCutFileName{\def\CommentCutFile{\jobname.cut}}
\usepackage{needspace}
\usepackage{xcolor}
\usepackage{microtype}
\usepackage{booktabs}

% ====================== CONFIGURATION ======================
\newif\ifsolucoes
\solucoesfalse                 % <<<< \solucoestrue for the ANSWER KEY; \solucoesfalse for the LIST

\newcommand{\disciplina}{Econometria Avançada --- MPE (Insper)}
\newcommand{\listatema}{L08 --- Univariate time series (Class 8): textbook exercises}
\newcommand{\espacoresol}{4.5cm}   % height of the solving space in the blank list

% ====================== MACHINERY ======================
\ifsolucoes\includecomment{sol}\else\excludecomment{sol}\fi

\pagestyle{fancy}\fancyhf{}
\lhead{\footnotesize\disciplina}
\rhead{\footnotesize List L08\ifsolucoes{} --- Solutions\fi}
\cfoot{\footnotesize\thepage}
\renewcommand{\headrulewidth}{0.3pt}
\setlength{\headheight}{14pt}

\setlist{topsep=2pt,itemsep=2pt}
\newcounter{questao}
% Equations are numbered per question, (Q<question>.<n>), so they never collide with the
% books' own numbers, which keep their printed tags, e.g. (3.7) or (18.21).
\renewcommand{\theequation}{Q\thequestao.\arabic{equation}}

% Block header
\newcommand{\bloco}[3]{\bigskip\noindent
  {\large\bfseries\color{blue!55!black}Block #1 --- #2}\hfill\textit{(#3)}\par
  \noindent\rule{\linewidth}{0.3pt}\smallskip}

% Question (automatic numbering; resets the equation counter)
\newcommand{\questao}[1]{\refstepcounter{questao}\setcounter{equation}{0}\medskip\needspace{5\baselineskip}\noindent
  \textbf{Question \thequestao.}~#1\par}
% Source line that opens every question (estilo-avaliacao §5, book questions)
\newcommand{\fonte}[1]{\par\noindent{\small\textit{Source:} #1}\par\smallskip}

% Notation list that opens every solution (estilo-avaliacao §4, "Equations, notation and layout").
\newenvironment{notacao}{\par\smallskip\noindent\textbf{Notation.}\begin{itemize}[leftmargin=1.4em,topsep=1pt,itemsep=0pt]\small}{\end{itemize}}
\allowdisplaybreaks

% Solving space in the blank list
\newcommand{\resolespaco}{\ifsolucoes\else\par\textit{Solution:}\vspace{\espacoresol}\par\fi}

% Solution step (use inside sol)
\newcommand{\passo}[2]{\par\smallskip\noindent\textbf{Step --- #1.}~#2}
\newcommand{\intuicao}[1]{\par\smallskip\noindent\textit{Intuition.}~#1}
\newcommand{\refbib}[1]{\par\smallskip\noindent{\footnotesize\color{black!60}\textbf{Ref:} #1}}

% =============================================================================
\begin{document}
\begin{center}
  {\Large\bfseries \disciplina}\\[2pt]
  {\large Exercise list: \listatema}\ifsolucoes\\[2pt]\textit{Solutions}\fi
\end{center}
\vspace{2pt}

{\footnotesize\noindent\textbf{Instructions.} These are the \emph{Practice} items of Class 8 in
\texttt{Map/class-plan.md} (T 2.1--2.6; W 18.C13; T 3.3, 3.4, 3.6, 3.11, 3.12, 3.14; W 12.4), grouped by
book. Every question opens with a source line, and the statement is reproduced verbatim from the book PDF.
Data exercises use Tsay's data files (FTS 3e web page, retrieved from the Internet Archive) and the
\texttt{wooldridge} package; see \texttt{Listas/data/README.md}. They are solved in Python
(\texttt{statsmodels}; GARCH-type models by conditional maximum likelihood written in NumPy/SciPy);
\texttt{Listas/L08-book-exercises-check.py} reproduces and asserts every number in the solutions. Tests
use the 5\% level unless the statement says otherwise. Show your steps.\par}

% =============================================================================
\bloco{1}{Tsay, Chapter 2: linear time series}{T Ch.~2 \,\textbullet\, Questions 1--6}

% ---------------------------------------------------------------- Q1
\questao{Tsay, Exercise 2.1 (MA(1) forecasts).}
\fonte{Tsay (2010), \textit{Analysis of Financial Time Series}, 3rd ed., Chapter 2, Exercise 2.1, printed
p.~104, reproduced verbatim. The chapter's exercises open with: ``If not specifically specified, use 5\%
significance level to draw conclusions in the exercises.''}

\noindent\textbf{2.1.} Suppose that the simple return of a monthly bond index follows the MA(1) model
\[ R_t = a_t + 0.2a_{t-1}, \qquad \sigma_a = 0.025. \]
Assume that $a_{100} = 0.01$. Compute the 1-step- and 2-step-ahead forecasts of the return at the forecast
origin $t = 100$. What are the standard deviations of the associated forecast errors? Also compute the lag-1
and lag-2 autocorrelations of the return series.
\resolespaco
\begin{sol}
\begin{notacao}
\item $R_t$: monthly simple return; $a_t$: white-noise shock, mean $0$, variance $\sigma_a^2$, $\sigma_a=0.025$.
\item $\theta=0.2$: MA coefficient (written with a plus sign, as in the statement); $F_{100}$: information up to $t=100$.
\item $\hat R_{100}(\ell)=E(R_{100+\ell}\mid F_{100})$: $\ell$-step forecast; $e_{100}(\ell)=R_{100+\ell}-\hat R_{100}(\ell)$: forecast error.
\item $\gamma_\ell$, $\rho_\ell$: lag-$\ell$ autocovariance and autocorrelation of $R_t$.
\end{notacao}

\noindent The model is
\begin{equation}\label{q21:model} R_t = a_t + \theta a_{t-1}, \qquad \theta = 0.2 . \end{equation}

\passo{1-step forecast}{Writing \eqref{q21:model} at $t=101$,}
\begin{equation}\label{q21:R101} R_{101} = a_{101} + \theta a_{100}. \end{equation}
Taking $E(\cdot\mid F_{100})$ of \eqref{q21:R101}, with $E(a_{101}\mid F_{100})=0$ and $a_{100}$ known:
\begin{align}
\hat R_{100}(1) &= \theta a_{100} \label{q21:f1}\\
&= (0.2)(0.01) \notag\\
&= 0.002. \notag
\end{align}
\passo{1-step error}{Subtracting \eqref{q21:f1} from \eqref{q21:R101} gives $e_{100}(1)=a_{101}$, so}
\begin{align*}
\operatorname{sd}[e_{100}(1)] &= \sigma_a\\
&= 0.025.
\end{align*}
\passo{2-step forecast}{Writing \eqref{q21:model} at $t=102$,}
\begin{equation}\label{q21:R102} R_{102} = a_{102} + \theta a_{101}. \end{equation}
Both shocks in \eqref{q21:R102} are future at $t=100$, so $\hat R_{100}(2)=0$ and $e_{100}(2)=a_{102}+\theta a_{101}$.
Because $a_{102}$ and $a_{101}$ are uncorrelated,
\begin{align}
\operatorname{Var}[e_{100}(2)] &= (1+\theta^2)\sigma_a^2, \label{q21:v2}\\
\operatorname{sd}[e_{100}(2)] &= \sigma_a\sqrt{1+\theta^2} \notag\\
&= (0.025)\sqrt{1+(0.2)^2} \notag\\
&= (0.025)\sqrt{1.04} \notag\\
&= 0.025495. \notag
\end{align}
\passo{autocorrelations}{From \eqref{q21:model}, $\gamma_0=(1+\theta^2)\sigma_a^2$ (as in \eqref{q21:v2}),
$\gamma_1=E[(a_t+\theta a_{t-1})(a_{t-1}+\theta a_{t-2})]=\theta\sigma_a^2$ and $\gamma_2=0$ (no shock in common). Dividing,}
\begin{align*}
\rho_1 &= \frac{\theta}{1+\theta^2}\\
&= \frac{0.2}{1+(0.2)^2}\\
&= \frac{0.2}{1.04}\\
&= 0.19231,\\
\rho_2 &= \frac{\gamma_2}{\gamma_0} = 0.
\end{align*}
\intuicao{An MA(1) has a memory of exactly one period: the shock $a_{100}$ moves only the next return, so
from two steps on the forecast reverts to the mean (zero) and the error variance reaches its ceiling
$\gamma_0$. The ACF cuts off after lag 1 for the same reason.}
\refbib{Tsay (2010) \S2.5.1 (p.~59, ACF of MA models) and \S2.5.4 (pp.~62--63, forecasting with MA models).}
\end{sol}

% ---------------------------------------------------------------- Q2
\questao{Tsay, Exercise 2.2 (AR(2) with $\phi_1=0$: moments and forecasts).}
\fonte{Tsay (2010), \textit{Analysis of Financial Time Series}, 3rd ed., Chapter 2, Exercise 2.2, printed
p.~104, reproduced verbatim.}

\noindent\textbf{2.2.} Suppose that the daily log return of a security follows the model
\[ r_t = 0.01 + 0.2r_{t-2} + a_t, \]
where $\{a_t\}$ is a Gaussian white noise series with mean zero and variance 0.02. What are the mean and
variance of the return series $r_t$? Compute the lag-1 and lag-2 autocorrelations of $r_t$. Assume that
$r_{100} = -0.01$, and $r_{99} = 0.02$. Compute the 1- and 2-step-ahead forecasts of the return series at the
forecast origin $t = 100$. What are the associated standard deviations of the forecast errors?
\resolespaco
\begin{sol}
\begin{notacao}
\item $r_t$: daily log return; $\phi_0=0.01$: intercept; $\phi_1=0$, $\phi_2=0.2$: AR coefficients (the model is an AR(2) without lag 1).
\item $a_t$: Gaussian white noise, $\sigma_a^2=0.02$; $\mu=E(r_t)$; $\gamma_\ell$, $\rho_\ell$: autocovariance and autocorrelation.
\item $\hat r_{100}(\ell)$: $\ell$-step forecast at origin 100; $e_{100}(\ell)$: its error.
\end{notacao}
\begin{equation}\label{q22:model} r_t = \phi_0 + \phi_2 r_{t-2} + a_t, \qquad \phi_0=0.01,\ \phi_2=0.2,\ \sigma_a^2=0.02. \end{equation}
The characteristic equation $1-0.2x^2=0$ has roots $\pm\sqrt{5}$, outside the unit circle, so the process is
weakly stationary and the moments below exist.

\passo{mean}{Taking expectations of \eqref{q22:model} with $E(r_t)=E(r_{t-2})=\mu$:}
\begin{align}
\mu &= \phi_0 + \phi_2\mu \label{q22:mu0}\\
\mu &= \frac{\phi_0}{1-\phi_2} \notag\\
&= \frac{0.01}{1-0.2} \notag\\
&= 0.0125. \notag
\end{align}
\passo{variance}{Taking the variance of \eqref{q22:model}; $a_t$ is uncorrelated with $r_{t-2}$ and $\operatorname{Var}(r_t)=\operatorname{Var}(r_{t-2})=\gamma_0$:}
\begin{align}
\gamma_0 &= \phi_2^2\gamma_0 + \sigma_a^2 \label{q22:g0}\\
\gamma_0 &= \frac{\sigma_a^2}{1-\phi_2^2} \notag\\
&= \frac{0.02}{1-(0.2)^2} \notag\\
&= \frac{0.02}{0.96} \notag\\
&= 0.020833. \notag
\end{align}
\passo{ACF}{Multiplying the mean-adjusted form of \eqref{q22:model} by $(r_{t-\ell}-\mu)$ and taking
expectations gives, for $\ell>0$, the Yule--Walker recursion with $\phi_1=0$:}
\begin{equation}\label{q22:yw} \rho_\ell = \phi_2\rho_{\ell-2}. \end{equation}
At $\ell=1$, \eqref{q22:yw} with $\rho_{-1}=\rho_1$ gives
\begin{align*}
\rho_1 &= \phi_2\rho_1\\
\rho_1(1-0.2) &= 0\\
\rho_1 &= 0.
\end{align*}
At $\ell=2$, \eqref{q22:yw} with $\rho_0=1$ gives
\begin{align*}
\rho_2 &= (0.2)(1)\\
&= 0.2.
\end{align*}
\passo{1-step forecast}{Writing \eqref{q22:model} at $t=101$ and conditioning on $F_{100}$ (here $r_{99}$ is known):}
\begin{align}
\hat r_{100}(1) &= \phi_0 + \phi_2 r_{99} \label{q22:f1}\\
&= 0.01 + 0.2(0.02) \notag\\
&= 0.01 + 0.004 \notag\\
&= 0.014. \notag
\end{align}
\passo{2-step forecast}{At $t=102$ the lag-2 value is $r_{100}$, also known:}
\begin{align}
\hat r_{100}(2) &= \phi_0 + \phi_2 r_{100} \label{q22:f2}\\
&= 0.01 + 0.2(-0.01) \notag\\
&= 0.01 - 0.002 \notag\\
&= 0.008. \notag
\end{align}
\passo{forecast errors}{Subtracting \eqref{q22:f1} and \eqref{q22:f2} from \eqref{q22:model} at $t=101,102$
leaves $e_{100}(1)=a_{101}$ and $e_{100}(2)=a_{102}$ (the $\psi_1=\phi_1$ weight is zero). Hence}
\begin{align*}
\operatorname{sd}[e_{100}(1)] = \operatorname{sd}[e_{100}(2)] &= \sqrt{0.02}\\
&= 0.141421.
\end{align*}
\intuicao{With only a lag-2 term the series is two interleaved AR(1)s (odd and even days). That is why
$\rho_1=0$ while $\rho_2=0.2$, and why both 1- and 2-step forecasts use an observed value and carry only
one shock of error.}
\refbib{Tsay (2010) \S2.4.1 (pp.~38--45, moments and ACF of AR models) and \S2.4.4 (pp.~54--56, AR forecasts).}
\end{sol}

% ---------------------------------------------------------------- Q3
\questao{Tsay, Exercise 2.3 (U.S. unemployment rate: model, forecasts, business cycles).}
\fonte{Tsay (2010), \textit{Analysis of Financial Time Series}, 3rd ed., Chapter 2, Exercise 2.3, printed
p.~104, reproduced verbatim.}

\noindent\textbf{2.3.} Consider the monthly U.S. unemployment rate from January 1948 to March 2009 in the
file \texttt{m-unrate.txt}. The data are seasonally adjusted and obtained from the Federal Reserve Bank of
St Louis. Build a time series model for the series and use the model to forecast the unemployment rate for
the April, May, June, and July of 2009. In addition, does the fitted model imply the existence of business
cycles? Why? (Note that there are more than one model fits the data well. You only need an adequate model.)
\resolespaco
\begin{sol}
\begin{notacao}
\item $x_t$: unemployment rate (\%), $t=1,\dots,735$ (1948:01--2009:03); $B$: back-shift operator, $Bx_t=x_{t-1}$.
\item $\phi_1,\phi_2$: nonseasonal AR coefficients of $(1-B)x_t$; $\theta_1,\theta_2$: nonseasonal MA coefficients;
$\Phi_1$, $\Theta_1$: seasonal (lag-12) AR and MA coefficients; $a_t$: white noise, variance $\sigma_a^2$.
\item $\omega_1,\omega_2$: characteristic (inverse) roots of the AR(2) factor; $\Delta=\phi_1^2+4\phi_2$: its discriminant.
\item $Q(m)$: Ljung--Box statistic on residuals with $m$ lags.
\end{notacao}

\passo{identification}{The ADF test (constant, 12 lags chosen by AIC) on the level gives $t=-2.783$,
$p=0.061$: a unit root is not rejected at 5\%, so we model $(1-B)x_t$. The residuals of nonseasonal models
keep a significant lag-12 autocorrelation (the seasonal adjustment is imperfect), so a seasonal $(1,0,1)_{12}$
factor is added. Among ARIMA$(p,1,q)\times(P,0,Q)_{12}$ models with $p\le 3$, $q\le 2$ we keep the one whose
residuals pass both $Q(12)$ and $Q(24)$ at 5\%.}

\noindent\textit{Code:}
\begingroup\footnotesize
\begin{verbatim}
x = pd.read_csv("Listas/data/m-unrate.txt", sep=r"\s+").Rate.values
m = ARIMA(x, order=(2,1,2), seasonal_order=(1,0,1,12)).fit()
acorr_ljungbox(m.resid[1:], lags=[12,24], model_df=6)
m.get_forecast(4).predicted_mean, m.get_forecast(4).se_mean
\end{verbatim}
\endgroup
\noindent\textit{Key output:}
\begingroup\footnotesize
\begin{verbatim}
ar.L1 1.1701 (0.1231)  ar.L2 -0.3104 (0.1194)  ma.L1 -1.1809 (0.1193)  ma.L2 0.4923 (0.1002)
ar.S.L12 0.5505 (0.0673)  ma.S.L12 -0.8134 (0.0468)  sigma2 0.0369
Q(12) = 12.52 (p = 0.051)   Q(24) = 26.43 (p = 0.090)     [df = 6 and 18]
\end{verbatim}
\endgroup

\passo{fitted model}{With standard errors in parentheses above, the fitted model is}
\begin{equation}\label{q23:model}
(1-1.1701B+0.3104B^2)(1-0.5505B^{12})(1-B)x_t=(1-1.1809B+0.4923B^2)(1-0.8134B^{12})a_t ,
\end{equation}
with $\hat\sigma_a^2=0.0369$.
All coefficients are significant; $Q(12)$ and $Q(24)$ do not reject white-noise residuals at 5\%, so
\eqref{q23:model} is adequate.

\passo{forecasts}{At origin 2009:03 (last value $8.6$), the forecasts of \eqref{q23:model} and their standard errors are}
\begin{center}\small
\begin{tabular}{lcccc}
\toprule
 & Apr 2009 & May 2009 & Jun 2009 & Jul 2009\\
\midrule
Forecast (\%) & 8.844 & 9.023 & 9.151 & 9.257\\
Std.\ error & 0.192 & 0.270 & 0.350 & 0.437\\
\bottomrule
\end{tabular}
\end{center}

\passo{business cycles}{Stochastic cycles require a pair of complex characteristic roots in the AR
polynomial of the (differenced) series. The nonseasonal factor of \eqref{q23:model} is
$1-\phi_1B-\phi_2B^2$ with, to five decimals, $\phi_1=1.17009$ and $\phi_2=-0.31036$. Its roots are complex if and only if}
\begin{equation}\label{q23:disc} \Delta=\phi_1^2+4\phi_2<0 . \end{equation}
Substituting the estimates into \eqref{q23:disc}:
\begin{align*}
\Delta &= (1.17009)^2 + 4(-0.31036)\\
&= 1.36911 - 1.24144\\
&= 0.12767 > 0 .
\end{align*}
The roots are real. The characteristic roots are
\begin{align*}
\omega_{1,2} &= \frac{\phi_1\pm\sqrt{\Delta}}{2}\\
&= \frac{1.17009\pm0.35731}{2}\\
&= 0.7637,\ 0.4064 .
\end{align*}
\noindent\textbf{Answer:} no. The nonseasonal AR part is a product of two AR(1) factors whose ACF is a mix
of two exponential decays, with no damped sine wave, so the fitted model does not imply business cycles.
The seasonal factor $1-0.5505B^{12}$ does have complex roots, but they sit at the seasonal frequencies
(periods $12/k$ months). They reflect leftover seasonality in the adjusted data, not a business cycle.
\intuicao{Unemployment rises fast in recessions and falls slowly in expansions. A linear ARIMA
fitted to monthly changes picks up this persistence as two real roots (0.76 and 0.41) rather than as an
oscillation. Another adequate model (e.g.\ an AR(3) on $x_t$) could have complex roots, which is why the
statement warns that more than one model fits.}
\refbib{Tsay (2010) \S2.4.1 (pp.~42--43, complex characteristic roots and business cycles), \S2.7.5 (pp.~76--80, unit-root test), \S2.8.2 (pp.~84--90, multiplicative seasonal models).}
\end{sol}

% ---------------------------------------------------------------- Q4
\questao{Tsay, Exercise 2.4 (CRSP decile returns: Ljung--Box test, ARMA model, forecasts).}
\fonte{Tsay (2010), \textit{Analysis of Financial Time Series}, 3rd ed., Chapter 2, Exercise 2.4, printed
p.~105, reproduced verbatim. The data file (not named in the statement) is Tsay's \texttt{m-deciles08.txt}.}

\noindent\textbf{2.4.} Consider the monthly simple returns of the Decile 1, Decile 2, Decile 9, and Decile 10
of NYSE/AMEX/NASDAQ based on market capitalization. The data span is from January 1970 to December 2008, and
the data are obtained from CRSP.
\begin{enumerate}[label=(\alph*)]
\item For the return series of Decile 2 and Decile 10, test the null hypothesis that the first 12 lags of
autocorrelations are zero at the 5\% level. Draw your conclusion.
\item Build an ARMA model for the return series of Decile 2. Perform model checking and write down the fitted model.
\item Use the fitted ARMA model to produce 1- to 12-step-ahead forecasts of the series and the associated
standard errors of forecasts.
\end{enumerate}
\resolespaco
\begin{sol}
\begin{notacao}
\item $R_t$: monthly simple return of Decile 2 (\texttt{CAP2RET}), $T=468$ (1970:01--2008:12).
\item $Jan_t$: January dummy; $\beta_0$, $\beta_1$: intercept and January effect; $u_t$: regression error, MA(1) with coefficient $\theta$.
\item $\hat\rho_\ell$: sample ACF; $Q(m)$: Ljung--Box statistic, $\chi^2_{m-g}$ under $H_0$ with $g$ fitted ARMA coefficients.
\end{notacao}

\passo{(a) Ljung--Box}{The test statistic is}
\begin{equation}\label{q24:Q} Q(m)=T(T+2)\sum_{\ell=1}^{m}\frac{\hat\rho_\ell^2}{T-\ell}\ \overset{H_0}{\sim}\ \chi^2_m . \end{equation}
Computing \eqref{q24:Q} with $m=12$ ($\chi^2_{12}$ 5\% critical value $21.03$):
\begin{center}\small
\begin{tabular}{lcc}
\toprule
 & $Q(12)$ & $p$-value\\
\midrule
Decile 2 & 55.74 & $1.3\times10^{-7}$\\
Decile 10 & 10.69 & 0.556\\
\bottomrule
\end{tabular}
\end{center}
We reject $H_0$ for Decile 2 (small caps have serial correlation) and do not reject for Decile 10 (largest
stocks: returns are serially uncorrelated, as in an efficient market for liquid stocks).

\passo{(b) identification}{For Decile 2, $\hat\rho_1=0.204$, $\hat\rho_{12}=0.231$ and $\hat\rho_{24}=0.130$
all exceed $2/\sqrt{T}=0.092$: there is a lag-1 and a seasonal (lag-12) pattern. A seasonal ARMA, as in Tsay's
Example 2.4 for Decile 1, gives}
\begin{equation}\label{q24:sar} (1-0.1981B)(1-0.9894B^{12})(R_t-0.0095)=(1-0.9398B^{12})a_t . \end{equation}
The seasonal AR and MA factors of \eqref{q24:sar} nearly cancel ($0.9894$ vs $0.9398$). That points to a
deterministic seasonal, the January effect, so we replace the seasonal factors by a January dummy and keep
an MA(1) for the lag-1 correlation (MA(1) has a lower BIC than AR(1) here):
\begin{equation}\label{q24:reg} R_t=\beta_0+\beta_1 Jan_t+u_t,\qquad u_t=a_t+\theta a_{t-1}. \end{equation}

\noindent\textit{Code:}
\begingroup\footnotesize
\begin{verbatim}
d   = pd.read_csv("Listas/data/m-deciles08.txt", sep=r"\s+")
jan = (d.date.astype(str).str[4:6] == "01").astype(float).values
m   = ARIMA(d.CAP2RET.values, exog=jan, order=(0,0,1), trend="c").fit()
acorr_ljungbox(m.resid, lags=[12,24], model_df=1)
ex = np.zeros(12); ex[0] = 1                   # step 1 = January 2009
f  = m.get_forecast(12, exog=ex)
\end{verbatim}
\endgroup
\noindent\textit{Key output:}
\begingroup\footnotesize
\begin{verbatim}
const 0.00388 (0.00395)   Jan 0.07891 (0.00789)   ma.L1 0.23304 (0.03681)   sigma2 0.00356
Q(12) = 14.50 (p = 0.207)   Q(24) = 23.65 (p = 0.424)      last residual a_T = -0.00406
\end{verbatim}
\endgroup
\passo{(b) fitted model and checking}{The fitted model \eqref{q24:reg} is}
\begin{equation}\label{q24:fit} R_t = 0.00388 + 0.07891\,Jan_t + a_t + 0.23304\,a_{t-1},\qquad \hat\sigma_a=0.0597 . \end{equation}
The January effect is about 7.9 percentage points ($t\approx10$). The residual $Q(12)=14.50$ ($p=0.207$) and
$Q(24)=23.65$ ($p=0.424$) do not reject white noise: the model is adequate.

\passo{(c) forecasts}{The origin is December 2008, so step 1 is January 2009. From \eqref{q24:fit}, the
1-step forecast uses the last residual $a_T=-0.00406$:}
\begin{align*}
\hat R_T(1) &= \beta_0+\beta_1+\theta a_T\\
&= 0.00388+0.07891+0.23304(-0.00406)\\
&= 0.08184 .
\end{align*}
For $\ell=2,\dots,12$ (February--December 2009) $Jan=0$ and the MA(1) memory is gone, so $\hat R_T(\ell)=\beta_0=0.00388$.
The standard errors are $\sigma_a=0.05965$ at $\ell=1$ and $\sigma_a\sqrt{1+\theta^2}$ for $\ell\ge2$:
\begin{align*}
\operatorname{se}(\ell\ge2) &= 0.05965\sqrt{1+(0.23304)^2}\\
&= 0.06125 .
\end{align*}
\begin{center}\small
\begin{tabular}{lcccc}
\toprule
Step $\ell$ & 1 (Jan) & 2 (Feb) & $\cdots$ & 12 (Dec)\\
\midrule
Forecast & 0.08184 & 0.00388 & $\cdots$ & 0.00388\\
Std.\ error & 0.05965 & 0.06125 & $\cdots$ & 0.06125\\
\bottomrule
\end{tabular}
\end{center}
\intuicao{Small-cap returns are predictable mainly through the calendar: the turn-of-the-year (January)
effect. When a seasonal AR factor near 1 is almost cancelled by a seasonal MA factor, the seasonality is
deterministic, and a dummy captures it with one parameter instead of two near-unit roots.}
\refbib{Tsay (2010) \S2.2 (pp.~32--33, Ljung--Box test), \S2.8.2 Example 2.4 (pp.~86--88, Decile 1 and the January effect), \S2.9 (pp.~90--97, regression with time series errors).}
\end{sol}

% ---------------------------------------------------------------- Q5
\questao{Tsay, Exercise 2.5 (long-range dependence in $|R_t|$ of IBM).}
\fonte{Tsay (2010), \textit{Analysis of Financial Time Series}, 3rd ed., Chapter 2, Exercise 2.5, printed
p.~105, reproduced verbatim.}

\noindent\textbf{2.5.} Consider the daily simple returns of IBM stock from 1970 to 2008 in the file
\texttt{d-ibm3dx7008.txt}. Compute the first 100 lags of ACF of the absolute series of daily simple returns of
IBM stock. Is there evidence of long-range dependence? Why?
\resolespaco
\begin{sol}
\begin{notacao}
\item $R_t$: IBM daily simple return, $T=9845$ (1970--2008); $|R_t|$: its absolute value.
\item $\hat\rho_k$: sample ACF of $|R_t|$ at lag $k$; $2/\sqrt T$: approximate 95\% band under white noise.
\item $d$: fractional-differencing parameter of a long-memory process, $-0.5<d<0.5$.
\end{notacao}

\noindent\textit{Code:}
\begingroup\footnotesize
\begin{verbatim}
r   = pd.read_csv("Listas/data/d-ibm3dx7008.txt", sep=r"\s+").rtn.values
rho = acf(np.abs(r), nlags=100)
b, _ = np.polyfit(np.log(np.arange(1,101)), np.log(rho[1:]), 1)   # rho_k ~ k^(2d-1)
\end{verbatim}
\endgroup
\noindent\textit{Key output:}
\begin{center}\small
\begin{tabular}{lccccccc}
\toprule
Lag $k$ & 1 & 2 & 5 & 10 & 20 & 50 & 100\\
\midrule
$\hat\rho_k$ & 0.1947 & 0.1638 & 0.1723 & 0.1311 & 0.1379 & 0.1089 & 0.0786\\
\bottomrule
\end{tabular}
\end{center}
All 100 sample autocorrelations are positive and above the band $2/\sqrt{9845}=0.0202$; the smallest is $0.0607$.

\passo{benchmark}{A short-memory (stationary ARMA) ACF decays exponentially. For example, an AR(1)
matching $\hat\rho_1$ would give $\rho_{100}=(0.1947)^{100}\approx 9\times10^{-72}$, while the observed
$\hat\rho_{100}=0.0786$.}
\passo{hyperbolic decay}{For a fractionally integrated process, Tsay's property 3 gives}
\begin{equation}\label{q25:hyp} \rho_k \approx C\,k^{2d-1}\quad (k\to\infty),\qquad \log\rho_k\approx \log C+(2d-1)\log k . \end{equation}
Regressing $\log\hat\rho_k$ on $\log k$ ($k=1,\dots,100$) gives slope $b=-0.2188$. Matching the slope with \eqref{q25:hyp}:
\begin{align*}
2d-1 &= -0.2188\\
d &= \frac{-0.2188+1}{2}\\
&= 0.3906 .
\end{align*}
This is a rough, descriptive estimate inside $(0,0.5)$, the long-memory range.

\noindent\textbf{Answer:} yes. The ACF of $|R_t|$ is small, positive and significant out to lag 100, and it
decays slowly (polynomially), not exponentially. That is the signature of long-range dependence in volatility,
even though the returns themselves are nearly uncorrelated.
\intuicao{The absolute return is a proxy for volatility. Volatility regimes last for months or years
(e.g.\ 1987, 2000--2002, 2008), so today's $|R_t|$ still says something about $|R_{t+100}|$. This is the
fact that motivates persistent GARCH and long-memory volatility models.}
\refbib{Tsay (2010) \S2.11 (pp.~101--103, long-memory models; property 3, $\rho_k\approx k^{2d-1}$, and the CRSP $|R_t|$ example, Figure 2.22).}
\end{sol}

% ---------------------------------------------------------------- Q6
\questao{Tsay, Exercise 2.6 (seasonal model for electricity demand).}
\fonte{Tsay (2010), \textit{Analysis of Financial Time Series}, 3rd ed., Chapter 2, Exercise 2.6, printed
p.~105, reproduced verbatim.}

\noindent\textbf{2.6.} Consider the demand of electricity of a manufacturing sector in the United States. The
data are logged, denote the demand of a fixed day of each month, and are in \texttt{power6.txt}. Build a time
series model for the series and use the fitted model to produce 1- to 24-step-ahead forecasts.
\resolespaco
\begin{sol}
\begin{notacao}
\item $x_t$: log electricity demand, monthly, $T=264$; $B$: back-shift operator; $s=12$: seasonal period.
\item $w_t=(1-B)(1-B^{12})x_t$: regular and seasonal difference.
\item $\theta$, $\Theta$: regular and seasonal MA coefficients (Tsay's sign convention $1-\theta B$); $a_t$: white noise with variance $\sigma_a^2$.
\end{notacao}

\passo{identification}{$x_t$ has a strong annual pattern: the ACF of $(1-B)x_t$ is $0.757$ at lag 12 and
$0.726$ at lag 24 and decays slowly at the seasonal lags. We therefore take $w_t=(1-B)(1-B^{12})x_t$. Its
ACF is significant only at lag 1 ($-0.332$), lag 12 ($-0.484$) and lag 13 ($0.206$), against a band of
$2/\sqrt{251}=0.126$. That is the pattern of the airline model:}
\begin{equation}\label{q26:air} (1-B)(1-B^{12})x_t=(1-\theta B)(1-\Theta B^{12})a_t , \end{equation}
whose $w_t$ has nonzero ACF only at lags 1, 11, 12 and 13, with $\rho_{13}=\rho_1\rho_{12}$
($(-0.332)(-0.484)=0.161$, close to the observed $0.206$).

\noindent\textit{Code:}
\begingroup\footnotesize
\begin{verbatim}
x = np.loadtxt("Listas/data/power6.txt")
m = ARIMA(x, order=(0,1,1), seasonal_order=(0,1,1,12)).fit()
acorr_ljungbox(m.resid[13:], lags=[12,24], model_df=2)
f = m.get_forecast(24); f.predicted_mean, f.se_mean
\end{verbatim}
\endgroup
\noindent\textit{Key output:}
\begingroup\footnotesize
\begin{verbatim}
ma.L1 -0.4874 (0.0574)   ma.S.L12 -0.9572 (0.1301)   sigma2 0.000337
Q(12) = 12.85 (p = 0.232)   Q(24) = 26.10 (p = 0.247)     [df = 10 and 22]
\end{verbatim}
\endgroup
\passo{fitted model}{In the notation of \eqref{q26:air}, $\hat\theta=0.4874$ and $\hat\Theta=0.9572$:}
\begin{equation}\label{q26:fit} (1-B)(1-B^{12})x_t=(1-0.4874B)(1-0.9572B^{12})a_t,\qquad \hat\sigma_a=0.0184 . \end{equation}
Both coefficients are significant, and $Q(12)$, $Q(24)$ do not reject white-noise residuals, so the model is
adequate. $\hat\Theta$ close to 1 suggests the seasonal pattern is close to deterministic: the seasonal
difference is nearly cancelled by the seasonal MA, the same remark Tsay makes about the airline model.

\passo{forecast recursion}{Expanding \eqref{q26:fit} and taking $E(\cdot\mid F_T)$, the 1-step forecast is}
\begin{equation}\label{q26:f1} \hat x_T(1)=x_T+x_{T-11}-x_{T-12}-\theta a_T-\Theta a_{T-11}+\theta\Theta a_{T-12}, \end{equation}
and for larger $\ell$ the same recursion is used with future $x$ replaced by forecasts and future $a$ by zero.

\passo{forecasts}{1- to 24-step-ahead forecasts of log demand, with standard errors:}
\begin{center}\small
\begin{tabular}{rcc@{\qquad}rcc}
\toprule
$\ell$ & $\hat x_T(\ell)$ & s.e. & $\ell$ & $\hat x_T(\ell)$ & s.e.\\
\midrule
1 & 9.8827 & 0.0185 & 13 & 9.9115 & 0.0380\\
2 & 9.8675 & 0.0208 & 14 & 9.8963 & 0.0392\\
3 & 9.8740 & 0.0228 & 15 & 9.9029 & 0.0405\\
4 & 9.8931 & 0.0247 & 16 & 9.9219 & 0.0417\\
5 & 9.9221 & 0.0265 & 17 & 9.9509 & 0.0429\\
6 & 10.0066 & 0.0281 & 18 & 10.0355 & 0.0441\\
7 & 10.0465 & 0.0297 & 19 & 10.0753 & 0.0452\\
8 & 10.0635 & 0.0311 & 20 & 10.0923 & 0.0463\\
9 & 10.0729 & 0.0326 & 21 & 10.1017 & 0.0473\\
10 & 10.0070 & 0.0339 & 22 & 10.0358 & 0.0484\\
11 & 9.9520 & 0.0352 & 23 & 9.9808 & 0.0494\\
12 & 9.9291 & 0.0365 & 24 & 9.9579 & 0.0504\\
\bottomrule
\end{tabular}
\end{center}
The second year repeats the first year's seasonal profile shifted up by about $0.029$ (e.g.\ $9.9115-9.8827$),
the drift implied by the double difference. The standard errors grow, with a jump after $\ell=12$ when
the seasonal shock enters.
\intuicao{The airline model says: this month is last month plus last year's month-to-month change, with
two smoothing (MA) corrections. The forecasts therefore copy the most recent seasonal shape and carry the
current level forward.}
\refbib{Tsay (2010) \S2.8.1--2.8.2 (pp.~82--90, seasonal differencing, the airline model, Example 2.3).}
\end{sol}

% =============================================================================
\bloco{2}{Tsay, Chapter 3: conditional heteroscedastic models}{T Ch.~3 \,\textbullet\, Questions 7--12}

% ---------------------------------------------------------------- Q7
\questao{Tsay, Exercise 3.3 (Gaussian AR(1)--GARCH(1,1) likelihood).}
\fonte{Tsay (2010), \textit{Analysis of Financial Time Series}, 3rd ed., Chapter 3, Exercise 3.3, printed
p.~168, reproduced verbatim.}

\noindent\textbf{3.3.} Suppose that $r_1,\dots,r_n$ are observations of a return series that follows the
AR(1)-GARCH(1,1) model
\[ r_t=\mu+\phi_1r_{t-1}+a_t,\qquad a_t=\sigma_t\epsilon_t,\qquad \sigma_t^2=\alpha_0+\alpha_1a_{t-1}^2+\beta_1\sigma_{t-1}^2, \]
where $\epsilon_t$ is a standard Gaussian white noise series. Derive the conditional log-likelihood function of the data.
\resolespaco
\begin{sol}
\begin{notacao}
\item $r_t$: return, $t=1,\dots,n$; $\mu$: intercept; $\phi_1$: AR(1) coefficient, $|\phi_1|<1$.
\item $a_t$: shock; $\sigma_t^2$: conditional variance; $\epsilon_t\sim$ i.i.d.\ $N(0,1)$.
\item $\alpha_0>0$, $\alpha_1\ge0$, $\beta_1\ge0$, $\alpha_1+\beta_1<1$; $\boldsymbol\theta=(\mu,\phi_1,\alpha_0,\alpha_1,\beta_1)'$.
\item $F_{t-1}$: information up to $t-1$; $\ell(\boldsymbol\theta)$: conditional log-likelihood.
\end{notacao}
\begin{gather}
r_t=\mu+\phi_1r_{t-1}+a_t, \label{q33:mean}\\
a_t=\sigma_t\epsilon_t, \label{q33:a}\\
\sigma_t^2=\alpha_0+\alpha_1a_{t-1}^2+\beta_1\sigma_{t-1}^2. \label{q33:var}
\end{gather}

\passo{conditioning}{The first observation has no $r_0$, and \eqref{q33:var} needs a starting value.
We condition on $r_1$ and on a fixed $\sigma_1^2$ (e.g.\ the sample variance of the residuals). Given these,
\eqref{q33:mean} gives $a_t=r_t-\mu-\phi_1r_{t-1}$ for $t\ge2$, and \eqref{q33:var} then gives every $\sigma_t^2$
recursively, so both are known functions of $\boldsymbol\theta$ and the data up to $t-1$.}

\passo{conditional density}{Using \eqref{q33:a} with $\epsilon_t\sim N(0,1)$ and $\sigma_t$ known at $t-1$,}
\begin{equation}\label{q33:cond} r_t\mid F_{t-1}\ \sim\ N\!\left(\mu+\phi_1r_{t-1},\ \sigma_t^2\right),\qquad t=2,\dots,n . \end{equation}
\passo{factorization}{By the chain rule of probability,}
\begin{equation}\label{q33:chain} f(r_2,\dots,r_n\mid r_1,\sigma_1^2)=\prod_{t=2}^{n}f(r_t\mid F_{t-1}). \end{equation}
Substituting the normal density of \eqref{q33:cond} into \eqref{q33:chain}:
\begin{equation}\label{q33:L} f(r_2,\dots,r_n\mid r_1,\sigma_1^2)=\prod_{t=2}^{n}\frac{1}{\sqrt{2\pi\sigma_t^2}}\exp\!\left[-\frac{(r_t-\mu-\phi_1r_{t-1})^2}{2\sigma_t^2}\right]. \end{equation}
\passo{logs}{Taking logs of \eqref{q33:L}:}
\begin{equation}\label{q33:ll} \ell(\boldsymbol\theta)=-\frac12\sum_{t=2}^{n}\left[\ln(2\pi)+\ln\sigma_t^2+\frac{(r_t-\mu-\phi_1r_{t-1})^2}{\sigma_t^2}\right], \end{equation}
where, substituting $a_{t-1}$ from \eqref{q33:mean} into \eqref{q33:var},
\begin{equation}\label{q33:rec} \sigma_t^2=\alpha_0+\alpha_1\left(r_{t-1}-\mu-\phi_1r_{t-2}\right)^2+\beta_1\sigma_{t-1}^2,\qquad t\ge3, \end{equation}
and $\sigma_2^2=\alpha_0+\alpha_1a_1^2+\beta_1\sigma_1^2$ with the fixed starting values. Dropping the constant,
$\ell(\boldsymbol\theta)\propto-\frac12\sum_{t=2}^n\big[\ln\sigma_t^2+a_t^2/\sigma_t^2\big]$, which is maximized numerically.
The check script verifies \eqref{q33:ll} against the library normal log-density.
\intuicao{The returns are not independent, but each one is normal given its past, with a mean and a
variance that the past fixes. The likelihood is the product of these one-step-ahead densities: a
prediction-error decomposition. Days with a large predicted variance count for less in the squared term.}
\refbib{Tsay (2010) \S3.4.3 (pp.~120--121, ``Estimation'', Gaussian conditional likelihood) and \S3.5 (p.~133).}
\end{sol}

% ---------------------------------------------------------------- Q8
\questao{Tsay, Exercise 3.4 (Student-$t$ AR(1)--GARCH(1,1) likelihood).}
\fonte{Tsay (2010), \textit{Analysis of Financial Time Series}, 3rd ed., Chapter 3, Exercise 3.4, printed
p.~168, reproduced verbatim.}

\noindent\textbf{3.4.} In the equation in Exercise 3.3, assume that $\epsilon_t$ follows a standardized
Student-$t$ distribution with $v$ degrees of freedom. Derive the conditional log-likelihood function of the data.
\resolespaco
\begin{sol}
\begin{notacao}
\item As in Question 7, with $\epsilon_t$ i.i.d.\ standardized Student-$t$ with $v>2$ degrees of freedom (mean 0, variance 1).
\item $\Gamma(\cdot)$: gamma function; $\boldsymbol\theta=(\mu,\phi_1,\alpha_0,\alpha_1,\beta_1)'$; $f_v(\cdot)$: standardized-$t$ density.
\end{notacao}

\passo{standardized $t$}{If $x\sim t_v$, then $\epsilon=x/\sqrt{v/(v-2)}$ has unit variance, and its density (Tsay Eq.~(3.7)) is}
\begin{equation}\label{q34:ft} f_v(\epsilon)=\frac{\Gamma\!\left(\frac{v+1}{2}\right)}{\Gamma\!\left(\frac v2\right)\sqrt{(v-2)\pi}}\left(1+\frac{\epsilon^2}{v-2}\right)^{-\frac{v+1}{2}} . \end{equation}
\passo{change of variables}{By \eqref{q33:a}, $a_t=\sigma_t\epsilon_t$ with $\sigma_t$ known at $t-1$, so $f(a_t\mid F_{t-1})=f_v(a_t/\sigma_t)/\sigma_t$. Substituting $\epsilon=a_t/\sigma_t$ in \eqref{q34:ft}:}
\begin{equation}\label{q34:fa} f(r_t\mid F_{t-1})=\frac{\Gamma\!\left(\frac{v+1}{2}\right)}{\Gamma\!\left(\frac v2\right)\sqrt{(v-2)\pi}}\ \frac{1}{\sigma_t}\left(1+\frac{a_t^2}{(v-2)\sigma_t^2}\right)^{-\frac{v+1}{2}},\qquad a_t=r_t-\mu-\phi_1r_{t-1}. \end{equation}
\passo{log-likelihood}{Using the factorization \eqref{q33:chain} with \eqref{q34:fa} and taking logs:}
\begin{equation}\label{q34:ll}
\begin{aligned}
\ell(\boldsymbol\theta,v)=(n-1)\left[\ln\Gamma\!\left(\tfrac{v+1}{2}\right)-\ln\Gamma\!\left(\tfrac v2\right)-\tfrac12\ln\big((v-2)\pi\big)\right]\\
-\sum_{t=2}^{n}\left[\frac{v+1}{2}\ln\!\left(1+\frac{(r_t-\mu-\phi_1r_{t-1})^2}{(v-2)\sigma_t^2}\right)+\frac12\ln\sigma_t^2\right],
\end{aligned}
\end{equation}
with $\sigma_t^2$ from the recursion \eqref{q33:rec}. If $v$ is fixed in advance, the first line is a constant
and only the sum matters (Tsay Eq.~(3.8)). If $v$ is estimated jointly, the first line must stay, because it
depends on $v$. The check script verifies \eqref{q34:ll} against \texttt{scipy.stats.t} rescaled to unit variance.
\intuicao{The $t$ log-density penalizes a large standardized shock only logarithmically, not quadratically
as the normal does. Outliers therefore pull the GARCH parameters less, and $v$ measures how fat the tails
are: as $v\to\infty$, \eqref{q34:ll} tends to the Gaussian likelihood of Question 7.}
\refbib{Tsay (2010) \S3.4.3 (pp.~121--122, Eqs.~(3.7)--(3.8) and the joint likelihood with $v$).}
\end{sol}

% ---------------------------------------------------------------- Q9
\questao{Tsay, Exercise 3.6 (Merck monthly returns: serial correlation, ARCH effects, ARCH model).}
\fonte{Tsay (2010), \textit{Analysis of Financial Time Series}, 3rd ed., Chapter 3, Exercise 3.6, printed
p.~169, reproduced verbatim.}

\noindent\textbf{3.6.} The file \texttt{m-mrk4608.txt} contains monthly simple returns of Merck stock from
June 1946 to December 2008. The file has two columns denoting date and simple return. Transform the simple
returns to log returns.
\begin{enumerate}[label=(\alph*)]
\item Is there any evidence of serial correlations in the log returns? Use autocorrelations and 5\%
significance level to answer the question. If yes, remove the serial correlations.
\item Is there any evidence of ARCH effects in the log returns? Use the residual series if there are serial
correlations in part (a). Use Ljung--Box statistics for the squared returns (or residuals) with 6 and 12 lags
of autocorrelations and 5\% significance level to answer the question.
\item Identify an ARCH model for the data and fit the identified model. Write down the fitted model.
\end{enumerate}
\resolespaco
\begin{sol}
\begin{notacao}
\item $r_t=\ln(1+R_t)$: monthly log return of Merck, $T=751$ (1946:06--2008:12); $R_t$: simple return.
\item $\phi_1,\phi_8$: AR coefficients at lags 1 and 8; $a_t$: mean-equation shock; $\sigma_t^2$: conditional variance.
\item $\alpha_0,\alpha_1,\alpha_2,\alpha_3$: ARCH(3) coefficients; $\tilde a_t=a_t/\sigma_t$: standardized residual.
\end{notacao}

\noindent\textit{Code (the ARCH fit uses \texttt{garch\_fit} of the check script, conditional Gaussian MLE):}
\begingroup\footnotesize
\begin{verbatim}
r  = np.log1p(pd.read_csv("Listas/data/m-mrk4608.txt", sep=r"\s+").mrk.values)
acf(r, nlags=12); acorr_ljungbox(r, lags=[6,12])                          # (a)
m  = ARIMA(r, order=([1,8],0,0), trend="c").fit()                        # (a) remove
acorr_ljungbox(m.resid**2, lags=[6,12]); pacf(m.resid**2, nlags=12)      # (b), (c)
g  = garch_fit(r, ar_lags=[1,8], p=3, q=0)                               # (c) AR-ARCH(3)
\end{verbatim}
\endgroup

\passo{(a) serial correlation}{The ACF of $r_t$ exceeds the band $2/\sqrt{751}=0.073$ at lag 1
($-0.078$) and lag 8 ($0.104$). Ljung--Box: $Q(6)=10.19$ ($p=0.117$) and $Q(12)=27.24$ ($p=0.007$). So yes,
there is serial correlation at the 5\% level, driven by lags 1 and 8. To remove it we fit an AR with lags 1
and 8 only (statsmodels reports the mean $0.0108$):}
\begin{equation}\label{q36:ar} (1+0.0826B-0.1094B^{8})(r_t-0.0108)=a_t,\qquad \text{s.e.}(\hat\phi_1)=0.0339,\ \text{s.e.}(\hat\phi_8)=0.0347 . \end{equation}
The residuals of \eqref{q36:ar} give $Q(12)=14.03$ with $12-2=10$ df ($p=0.172$): no remaining serial correlation.

\passo{(b) ARCH effects}{On the squared residuals $a_t^2$ of \eqref{q36:ar}: $Q(6)=17.81$ ($p=0.0067$) and
$Q(12)=29.82$ ($p=0.0030$). Both reject $H_0$ of no autocorrelation in $a_t^2$: there are ARCH effects.}

\passo{(c) identification}{The PACF of $a_t^2$ at lags 1--4 is $0.056$, $0.056$, $0.110$, $0.045$, and
is inside the band after lag 4. Only lag 3 exceeds $0.073$, so we identify an ARCH(3) and estimate it jointly
with the AR mean equation:}
\begin{equation}\label{q36:arch} r_t=\mu+\phi_1r_{t-1}+\phi_8r_{t-8}+a_t,\quad a_t=\sigma_t\epsilon_t,\quad \sigma_t^2=\alpha_0+\alpha_1a_{t-1}^2+\alpha_2a_{t-2}^2+\alpha_3a_{t-3}^2 . \end{equation}
\noindent\textit{Key output (Gaussian conditional MLE; standard errors from the numerical Hessian):}
\begingroup\footnotesize
\begin{verbatim}
mu 0.01170 (0.00265)  phi1 -0.08134 (0.03706)  phi8 0.10145 (0.03715)
alpha0 0.00404 (0.00033)  alpha1 0.02741 (0.03856)  alpha2 0.05283 (0.03560)
alpha3 0.09608 (0.04142)
standardized residuals: Q(12) = 14.14 (p = 0.167)
squared:                Q(6) = 1.64 (p = 0.949), Q(12) = 11.20 (p = 0.512)
\end{verbatim}
\endgroup
\passo{fitted model}{Substituting the estimates into \eqref{q36:arch}:}
\begin{gather*}
r_t=0.0117-0.0813\,r_{t-1}+0.1015\,r_{t-8}+a_t,\qquad a_t=\sigma_t\epsilon_t,\ \epsilon_t\sim N(0,1),\\
\sigma_t^2=0.00404+0.0274\,a_{t-1}^2+0.0528\,a_{t-2}^2+0.0961\,a_{t-3}^2 .
\end{gather*}
The standardized residuals and their squares pass the Ljung--Box tests, so the model is adequate. The
implied unconditional variance is
\begin{align*}
\operatorname{Var}(a_t) &= \frac{\alpha_0}{1-\alpha_1-\alpha_2-\alpha_3}\\
&= \frac{0.00404}{1-(0.02741+0.05283+0.09608)}\\
&= \frac{0.00404}{0.82368}\\
&= 0.00490,
\end{align*}
a monthly volatility of about $7.0\%$. Only $\alpha_3$ is individually significant ($t=2.32$), so the ARCH
effect in monthly Merck returns is real but weak.
\intuicao{Monthly returns aggregate many daily shocks, which weakens volatility clustering. The ARCH
coefficients sum to only $0.18$, far from the persistence of 0.99 typical of daily data (Questions 10--12).}
\refbib{Tsay (2010) \S3.3.1 (p.~114, testing for ARCH effects), \S3.4.3 (pp.~119--123, building an ARCH model: PACF of $a_t^2$, estimation, model checking).}
\end{sol}

% ---------------------------------------------------------------- Q10
\questao{Tsay, Exercise 3.11 (GM daily log returns: ARCH test, PACF, normal and GED GARCH).}
\fonte{Tsay (2010), \textit{Analysis of Financial Time Series}, 3rd ed., Chapter 3, Exercise 3.11, printed
p.~170, reproduced verbatim.}

\noindent\textbf{3.11.} The file \texttt{d-gmsp9908.txt} contains the daily simple returns of GM stock and
the S\&P composite index from 1999 to 2008. It has three columns denoting date, GM return, and S\&P return.
\begin{enumerate}[label=(\alph*)]
\item Compute the daily log returns of GM stock. Is there any evidence of ARCH effects in the log returns?
You may use 10 lags of the squared returns and 5\% significance level to perform the test.
\item Compute the PACF of the squared log returns (10 lags).
\item Specify a GARCH model for the GM log return using a normal distribution for the innovations. Perform
model checking and write down the fitted model.
\item Find an adequate GARCH model for the series but using the generalized error distribution for the
innovations. Write down the fitted model.
\end{enumerate}
\resolespaco
\begin{sol}
\begin{notacao}
\item $r_t=100\ln(1+R_t)$: GM daily log return in percent, $T=2515$ (1999-01-04 to 2008-12-31). Percent
units keep $\alpha_0$ away from zero; only $\mu$ and $\alpha_0$ change scale.
\item $a_t=r_t-\mu$; $\sigma_t^2$: conditional variance; $\alpha_0,\alpha_1,\beta_1$: GARCH(1,1) coefficients.
\item $v$: GED shape parameter ($v=2$ is normal; $v<2$ has fat tails); $\tilde a_t=a_t/\sigma_t$.
\end{notacao}
\noindent\textit{Code:}
\begingroup\footnotesize
\begin{verbatim}
d  = pd.read_csv("Listas/data/d-gmsp9908.txt", sep=r"\s+")
r  = 100*np.log1p(d.gm.values); e = r - r.mean()
het_arch(e, nlags=10); acorr_ljungbox(e**2, lags=[10]); pacf(e**2, nlags=10)   # (a), (b)
gn = garch_fit(r)               # (c) GARCH(1,1), normal
gg = garch_fit(r, dist="ged")   # (d) GARCH(1,1), GED (Tsay Eq. (3.10))
\end{verbatim}
\endgroup

\passo{(a) ARCH test}{Engle's LM test regresses $e_t^2$ on 10 of its lags; under $H_0$ of no ARCH, $TR^2\sim\chi^2_{10}$.
We get $TR^2=625.2$ ($F=82.95$), with $p<10^{-100}$. The Ljung--Box statistic of $e_t^2$ is $Q(10)=898.1$.
Strong evidence of ARCH effects.}
\passo{(b) PACF of $e_t^2$}{(band $2/\sqrt{2515}=0.040$)}
\begin{center}\small
\begin{tabular}{lcccccccccc}
\toprule
Lag & 1 & 2 & 3 & 4 & 5 & 6 & 7 & 8 & 9 & 10\\
\midrule
PACF & 0.193 & 0.397 & $-0.013$ & $-0.128$ & 0.139 & 0.083 & 0.089 & 0.030 & $-0.000$ & 0.118\\
\bottomrule
\end{tabular}
\end{center}
Significant spikes out to lag 10 would require a long ARCH; a GARCH(1,1) is the parsimonious alternative.

\passo{(c) normal GARCH(1,1)}{The raw returns have $Q(10)=28.40$ ($p=0.0016$), but the standardized
residuals of a constant-mean GARCH(1,1) are white, so no AR term is needed. The model is}
\begin{equation}\label{q311:g} r_t=\mu+a_t,\qquad a_t=\sigma_t\epsilon_t,\qquad \sigma_t^2=\alpha_0+\alpha_1a_{t-1}^2+\beta_1\sigma_{t-1}^2 . \end{equation}
\noindent\textit{Key output:}
\begingroup\footnotesize
\begin{verbatim}
normal: mu 0.0245 (0.0437)  alpha0 0.0942 (0.0208)  alpha1 0.0701 (0.0091)
        beta1 0.9207 (0.0097)  logL -5853.30
  std. resid: Q(10) = 12.78 (p 0.236), Q(20) = 26.92 (p 0.137)
  squared:    Q(10) = 1.79 (p 0.998),  Q(20) = 9.01 (p 0.983)
  kurtosis of std. resid = 8.89
\end{verbatim}
\endgroup
Substituting into \eqref{q311:g}:
\begin{equation}\label{q311:gfit} r_t=0.0245+a_t,\qquad \sigma_t^2=0.0942+0.0701\,a_{t-1}^2+0.9207\,\sigma_{t-1}^2,\qquad \epsilon_t\sim N(0,1). \end{equation}
Model checking: the Ljung--Box tests on $\tilde a_t$ and $\tilde a_t^2$ do not reject, so the mean and
variance equations are adequate. The persistence $\hat\alpha_1+\hat\beta_1=0.9908$ is close to 1. The
standardized residuals have kurtosis $8.89$, far above the normal's $3$, so the normal assumption fails in the tails.

\passo{(d) GED GARCH(1,1)}{Same equations \eqref{q311:g}, with $\epsilon_t$ following the GED of Tsay Eq.~(3.10):}
\begingroup\footnotesize
\begin{verbatim}
GED:  mu -0.0444 (0.0408)  alpha0 0.0542 (0.0212)  alpha1 0.0678 (0.0113)
      beta1 0.9286 (0.0117)  v 1.2521 (0.0432)  logL -5747.55
  std. resid: Q(10) = 13.09 (p 0.219);  squared: Q(10) = 1.86 (p 0.997)
\end{verbatim}
\endgroup
\begin{equation}\label{q311:ged} r_t=-0.0444+a_t,\qquad \sigma_t^2=0.0542+0.0678\,a_{t-1}^2+0.9286\,\sigma_{t-1}^2,\qquad \epsilon_t\sim \text{GED}(v=1.252). \end{equation}
The test of normality inside the GED family, $H_0\!: v=2$, gives
\begin{align*}
t &= \frac{\hat v-2}{\text{s.e.}(\hat v)}\\
&= \frac{1.2521-2}{0.0432}\\
&= -17.33,
\end{align*}
a decisive rejection. The log-likelihood also rises by $105.75$ ($-5747.55$ vs $-5853.30$) for one extra
parameter. The residual checks pass, so \eqref{q311:ged} is adequate and preferred to \eqref{q311:gfit}.
\intuicao{GM's daily returns cluster in volatility (persistence near 0.99) and have fat tails even after
standardizing. The GED with $v\approx1.25$ (between Laplace, $v=1$, and normal, $v=2$) absorbs the tail
mass, and $\alpha_0$ falls because fewer extreme days have to be explained by the variance equation.}
\refbib{Tsay (2010) \S3.3.1 (p.~114, Engle's LM test), \S3.5 (pp.~131--135, GARCH), \S3.4.3 (p.~122, GED, Eq.~(3.10)).}
\end{sol}

% ---------------------------------------------------------------- Q11
\questao{Tsay, Exercise 3.12 (S\&P daily simple returns: ARCH test, GARCH model, forecasts).}
\fonte{Tsay (2010), \textit{Analysis of Financial Time Series}, 3rd ed., Chapter 3, Exercise 3.12, printed
p.~170, reproduced verbatim.}

\noindent\textbf{3.12.} Consider the daily simple returns of the S\&P composite index in the file
\texttt{d-gmsp9908.txt}.
\begin{enumerate}[label=(\alph*)]
\item Is there any ARCH effect in the simple return series? Use 10 lags of the squared returns and 5\%
significance level to perform the test.
\item Build an adequate GARCH model for the simple return series.
\item Compute 1-step- to 4-step-ahead forecasts of the simple return and its volatility based on the fitted model.
\end{enumerate}
\resolespaco
\begin{sol}
\begin{notacao}
\item $r_t=100R_t$: S\&P daily simple return in percent, $T=2515$, last date $T=$ 2008-12-31.
\item $\mu,\phi_1,\phi_2$: AR(2) mean equation; $a_t$: shock; $\sigma_t^2$: conditional variance; $\alpha_0,\alpha_1,\beta_1$: GARCH(1,1).
\item $\hat r_T(\ell)$, $\sigma_T^2(\ell)$: $\ell$-step forecasts of the return and of its conditional variance; volatility $=\sigma_T(\ell)$.
\end{notacao}
\noindent\textit{Code:}
\begingroup\footnotesize
\begin{verbatim}
r = 100*d.sp.values; e = r - r.mean()
het_arch(e, nlags=10); acorr_ljungbox(e**2, lags=[10])     # (a)
g = garch_fit(r, ar_lags=[1,2])                            # (b) AR(2)-GARCH(1,1), normal
\end{verbatim}
\endgroup
\passo{(a) ARCH test}{Engle's LM test with 10 lags: $TR^2=711.2$ ($F=98.88$), $p\approx0$; $Q(10)$ of the
squared demeaned returns $=2097.7$. Yes, strong ARCH effects.}
\passo{(b) model}{The returns have $\hat\rho_1=-0.077$, $\hat\rho_2=-0.091$ (band $0.040$) and $Q(10)=57.06$,
so the mean needs lags 1 and 2. The model is}
\begin{equation}\label{q312:m} r_t=\mu+\phi_1r_{t-1}+\phi_2r_{t-2}+a_t,\quad a_t=\sigma_t\epsilon_t,\quad \sigma_t^2=\alpha_0+\alpha_1a_{t-1}^2+\beta_1\sigma_{t-1}^2 . \end{equation}
\noindent\textit{Key output:}
\begingroup\footnotesize
\begin{verbatim}
mu 0.0357 (0.0179)  phi1 -0.0589 (0.0210)  phi2 -0.0419 (0.0208)
alpha0 0.0100 (0.0031)  alpha1 0.0714 (0.0094)  beta1 0.9235 (0.0101)
std. resid: Q(10) = 8.57 (p 0.380), Q(20) = 23.84 (p 0.160)   [df 8, 18]
squared:    Q(10) = 14.65 (p 0.145), Q(20) = 20.31 (p 0.438)
\end{verbatim}
\endgroup
Substituting into \eqref{q312:m}, with five decimals because they are used in the forecasts:
\begin{gather}
r_t=0.03570-0.05888\,r_{t-1}-0.04190\,r_{t-2}+a_t, \label{q312:mean}\\
\sigma_t^2=0.00997+0.07135\,a_{t-1}^2+0.92350\,\sigma_{t-1}^2,\qquad \epsilon_t\sim N(0,1). \label{q312:var}
\end{gather}
All Ljung--Box tests on $\tilde a_t$ and $\tilde a_t^2$ pass at 5\%: the model is adequate. Persistence
$\hat\alpha_1+\hat\beta_1=0.99485$, and the unconditional variance is $\alpha_0/(1-\alpha_1-\beta_1)=1.93$ (daily s.d.\ $\approx1.39\%$).

\passo{(c) inputs at the origin}{At $T$ (2008-12-31): $r_T=1.4158$, $r_{T-1}=2.4407$, $a_T=1.5076$, $\sigma_T^2=7.9263$.}
\passo{return forecasts}{Taking $E(\cdot\mid F_T)$ of \eqref{q312:mean}:}
\begin{align}
\hat r_T(1) &= \mu+\phi_1r_T+\phi_2r_{T-1} \label{q312:f1}\\
&= 0.03570+(-0.05888)(1.4158)+(-0.04190)(2.4407) \notag\\
&= 0.03570-0.08336-0.10227 \notag\\
&= -0.1499 . \notag
\end{align}
For $\ell\ge2$ the recursion is $\hat r_T(\ell)=\mu+\phi_1\hat r_T(\ell-1)+\phi_2\hat r_T(\ell-2)$:
\begin{align*}
\hat r_T(2) &= 0.03570+(-0.05888)(-0.1499)+(-0.04190)(1.4158) = -0.0148,\\
\hat r_T(3) &= 0.03570+(-0.05888)(-0.0148)+(-0.04190)(-0.1499) = 0.0429,\\
\hat r_T(4) &= 0.03570+(-0.05888)(0.0429)+(-0.04190)(-0.0148) = 0.0338 .
\end{align*}
\passo{volatility forecasts}{The variance recursion compounds rounding over the steps, so it is shown with six
decimals: $\alpha_0=0.009968$, $\alpha_1=0.071347$, $\beta_1=0.923498$, $a_T=1.507578$, $\sigma_T^2=7.926285$.
From \eqref{q312:var} at $T+1$, everything is known at $T$:}
\begin{align}
\sigma_T^2(1) &= \alpha_0+\alpha_1a_T^2+\beta_1\sigma_T^2 \label{q312:v1}\\
&= 0.009968+0.071347(1.507578)^2+0.923498(7.926285) \notag\\
&= 0.009968+0.162157+7.319908 \notag\\
&= 7.492033 . \notag
\end{align}
For $\ell\ge2$, $E(a_{T+\ell-1}^2\mid F_T)=\sigma_T^2(\ell-1)$, so
\begin{equation}\label{q312:vl} \sigma_T^2(\ell)=\alpha_0+(\alpha_1+\beta_1)\,\sigma_T^2(\ell-1),\qquad \alpha_1+\beta_1=0.994845 . \end{equation}
Substituting into \eqref{q312:vl}:
\begin{align*}
\sigma_T^2(2) &= 0.009968+0.994845(7.492033) = 7.463380,\\
\sigma_T^2(3) &= 0.009968+0.994845(7.463380) = 7.434874,\\
\sigma_T^2(4) &= 0.009968+0.994845(7.434874) = 7.406515 .
\end{align*}
\begin{center}\small
\begin{tabular}{lcccc}
\toprule
Step $\ell$ & 1 & 2 & 3 & 4\\
\midrule
Return forecast $\hat r_T(\ell)$ (\%) & $-0.1499$ & $-0.0148$ & $0.0429$ & $0.0338$\\
Variance $\sigma_T^2(\ell)$ & 7.4920 & 7.4634 & 7.4349 & 7.4065\\
Volatility $\sigma_T(\ell)$ (\%) & 2.7372 & 2.7319 & 2.7267 & 2.7215\\
\bottomrule
\end{tabular}
\end{center}
\intuicao{The origin is the end of 2008, at the peak of the crisis: the conditional variance (7.9) is about
four times its long-run level (1.93). With persistence $0.995$ the forecast moves back towards 1.93 only very
slowly, about $0.5\%$ of the gap per day, while the mean forecasts are tiny and quickly settle at
$\mu/(1-\phi_1-\phi_2)=0.0357/1.1008\approx0.032$.}
\refbib{Tsay (2010) \S3.5 (pp.~132--133, GARCH(1,1) forecasts, $\sigma_h^2(\ell)=\alpha_0+(\alpha_1+\beta_1)\sigma_h^2(\ell-1)$); \S2.4.4 (AR forecasts).}
\end{sol}

% ---------------------------------------------------------------- Q12
\questao{Tsay, Exercise 3.14 (market volatility as a regressor in GM's volatility).}
\fonte{Tsay (2010), \textit{Analysis of Financial Time Series}, 3rd ed., Chapter 3, Exercise 3.14, printed
p.~171, reproduced verbatim.}

\noindent\textbf{3.14.} Revisit the file \texttt{d-gmsp9908.txt}. However, we shall investigate the value of
using market volatility in modeling volatility of individual stocks. Convert the two simple return series into
percentage log return series.
\begin{enumerate}[label=(\alph*)]
\item Build an AR(5)--GARCH(1,1) model with generalized error distribution for the log S\&P returns. The
AR(5) contains only lags 3 and 5. Denote the fitted volatility series by \texttt{spvol}.
\item Estimate a GARCH(1,1) model with \texttt{spvol} as an exogenous variable to the log GM return series.
Check the adequacy of the model, and write down the fitted model. In S-Plus, the command is
\begin{center}\texttt{fit = garch(gm \textasciitilde 1, \textasciitilde garch(1,1)+spvol, cond.dist='ged')}\end{center}
\item Discuss the implication of the fitted model.
\end{enumerate}
\resolespaco
\begin{sol}
\begin{notacao}
\item $s_t$, $g_t$: S\&P and GM daily log returns in percent ($100\ln(1+R_t)$), 1999--2008.
\item $\sigma_{s,t}$: fitted S\&P conditional standard deviation = \texttt{spvol} (Tsay stores ``the volatility series $\sigma_t$'' in \texttt{sigma.t}).
\item $\sigma_{g,t}^2$: GM conditional variance; $\gamma$: coefficient of \texttt{spvol} in GM's variance equation; $v$: GED shape.
\end{notacao}
\noindent\textit{Code:}
\begingroup\footnotesize
\begin{verbatim}
s, g  = 100*np.log1p(d.sp.values), 100*np.log1p(d.gm.values)
fs    = garch_fit(s, ar_lags=[3,5], dist="ged")             # (a)
spvol = np.sqrt(fs["s2"])                                   # sigma_t, t = 6..T
fg    = garch_fit(g[5:], X=spvol[:,None], dist="ged")       # (b) gamma*spvol_t in the variance
f0    = garch_fit(g[5:], dist="ged")                        # same sample, gamma = 0
\end{verbatim}
\endgroup
\passo{(a) S\&P model}{The model is}
\begin{equation}\label{q314:sp} s_t=\mu+\phi_3s_{t-3}+\phi_5s_{t-5}+a_t,\quad a_t=\sigma_{s,t}\epsilon_t,\quad \sigma_{s,t}^2=\alpha_0+\alpha_1a_{t-1}^2+\beta_1\sigma_{s,t-1}^2,\quad \epsilon_t\sim\text{GED}(v) . \end{equation}
\noindent\textit{Key output:}
\begingroup\footnotesize
\begin{verbatim}
mu 0.0444 (0.0175)  phi3 -0.0094 (0.0207)  phi5 -0.0414 (0.0209)
alpha0 0.0079 (0.0033)  alpha1 0.0713 (0.0108)  beta1 0.9258 (0.0113)
v 1.4964 (0.0603)   logL -3688.80
std. resid: Q(10) = 12.80;  squared: Q(10) = 11.49
\end{verbatim}
\endgroup
\begin{gather*}
s_t=0.0444-0.0094\,s_{t-3}-0.0414\,s_{t-5}+a_t,\\
\sigma_{s,t}^2=0.0079+0.0713\,a_{t-1}^2+0.9258\,\sigma_{s,t-1}^2,\qquad v=1.496 .
\end{gather*}
The residual checks pass ($Q(10)=12.80$ with 8 df, $p=0.119$; $Q(10)=11.49$ for squares, $p=0.320$).
\texttt{spvol}$_t=\hat\sigma_{s,t}$, with mean $1.1749$.

\passo{(b) GM model with \texttt{spvol}}{The S-Plus formula puts \texttt{spvol} in the variance equation:}
\begin{equation}\label{q314:gm} g_t=\mu+a_t,\quad a_t=\sigma_{g,t}\epsilon_t,\quad \sigma_{g,t}^2=\alpha_0+\alpha_1a_{t-1}^2+\beta_1\sigma_{g,t-1}^2+\gamma\,\texttt{spvol}_t,\quad \epsilon_t\sim\text{GED}(v) . \end{equation}
\noindent\textit{Key output (sample $t=6,\dots,T$, where \texttt{spvol} exists):}
\begingroup\footnotesize
\begin{verbatim}
with spvol:  mu -0.0485 (0.0395)  alpha0 0.0214 (0.0355)  alpha1 0.0677 (0.0114)
             beta1 0.9254 (0.0127)  gamma 0.0499 (0.0474)  v 1.2478 (0.0430)   logL -5733.16
             std. resid: Q(10) = 12.44 (p 0.257);  squared: Q(10) = 1.87 (p 0.997)
gamma = 0:   alpha0 0.0540  alpha1 0.0677  beta1 0.9287  v 1.2505   logL -5733.79
\end{verbatim}
\endgroup
Substituting into \eqref{q314:gm}:
\begin{equation}\label{q314:fit} g_t=-0.0485+a_t,\qquad \sigma_{g,t}^2=0.0214+0.0677\,a_{t-1}^2+0.9254\,\sigma_{g,t-1}^2+0.0499\,\texttt{spvol}_t,\qquad v=1.248 . \end{equation}
The Ljung--Box tests on the standardized residuals and their squares do not reject, so \eqref{q314:fit} is
adequate. The Wald test of $H_0\!:\gamma=0$ is
\begin{align*}
t &= \frac{\hat\gamma}{\text{s.e.}(\hat\gamma)}\\
&= \frac{0.0499}{0.0474}\\
&= 1.05,
\end{align*}
and the likelihood-ratio statistic is
\begin{align*}
LR &= 2\,[\ell(\text{with spvol})-\ell(\gamma=0)]\\
&= 2\,[-5733.16-(-5733.79)]\\
&= 1.28 < \chi^2_{1,0.95}=3.84 .
\end{align*}
\passo{(c) implication}{\texttt{spvol} is not significant: once GM's own lagged shock $a_{t-1}^2$ and lagged
variance $\sigma_{g,t-1}^2$ are in the model, market volatility adds no information about GM's volatility.
Its average contribution $\hat\gamma\cdot\overline{\texttt{spvol}}=0.0499\times1.1749=0.0586$ mostly replaces
the intercept: $\hat\alpha_0$ falls from $0.0540$ to $0.0214$, while $\alpha_1$ and $\beta_1$ hardly move.
Market-wide volatility shocks also hit GM's returns, so GM's GARCH recursion already tracks them, and the
persistent \texttt{spvol} series acts almost like a constant in GM's variance equation.}
\intuicao{Volatility is largely a common factor (2008 raised every stock's volatility), but a univariate
GARCH filters that factor out of the stock's own squared shocks. To see a spillover, a regressor must carry
information that the stock's own past does not: here it does not.}
\refbib{Tsay (2010) \S3.5 (pp.~131--135, GARCH), \S3.4.3 (p.~122, GED, Eq.~(3.10)); for volatility regressors see the end-of-chapter S-Plus command quoted in the statement.}
\end{sol}

% =============================================================================
\bloco{3}{Wooldridge: unit roots and ARCH}{W Ch.~18 and 12 \,\textbullet\, Questions 13--14}

% ---------------------------------------------------------------- Q13
\questao{Wooldridge, Computer Exercise 18.C13 (Dickey--Fuller tests, TRAFFIC2).}
\fonte{Wooldridge (2016), \textit{Introductory Econometrics: A Modern Approach}, 6th ed., Chapter 18,
Computer Exercise C13, printed p.~604, reproduced verbatim.}

\noindent\textbf{C13} Use the data in TRAFFIC2 for this exercise. These monthly data, on traffic accidents in
California over the years 1981 to 1989, were used in Computer Exercise C11 in Chapter 10.
\begin{enumerate}[label=(\roman*)]
\item Using the standard Dickey-Fuller regression, test whether $ltotacc_t$ has a unit root. Can you reject a
unit root at the 2.5\% level?
\item Now, add two lagged changes to the test from part (i) and compute the augmented Dickey-Fuller test. What do you conclude?
\item Add a linear time trend to the ADF regression from part (ii). Now what happens?
\item Given the findings from parts (i) through (iii), what would you say is the best characterization of
$ltotacc_t$: an I(1) process or an I(0) process about a linear time trend?
\item Test the percentage of fatalities, $prcfat_t$, for a unit root, using two lags in an ADF regression. In
this case, does it matter whether you include a linear time trend?
\end{enumerate}
\resolespaco
\begin{sol}
\begin{notacao}
\item $y_t$: the series tested ($ltotacc_t$ = log of total accidents, or $prcfat_t$ = $100\times$ fatal/total accidents), $T=108$ months (1981--1989).
\item $\Delta y_t=y_t-y_{t-1}$; $\theta=\rho-1$: coefficient on $y_{t-1}$; $\gamma_j$: coefficients on lagged changes; $\delta$: trend coefficient.
\item $t_{\hat\theta}$: Dickey--Fuller $t$ statistic. Asymptotic critical values (Wooldridge Tables 18.2 and 18.3):
no trend $-3.43$ (1\%), $-3.12$ (2.5\%), $-2.86$ (5\%), $-2.57$ (10\%); with trend $-3.96$, $-3.66$, $-3.41$, $-3.12$.
\end{notacao}
The regressions, Wooldridge Eqs.~(18.21) and (18.24), plus a trend in (iii) and (v):
\begin{equation}\label{q13:adf} \Delta y_t=\alpha+\delta t+\theta y_{t-1}+\gamma_1\Delta y_{t-1}+\gamma_2\Delta y_{t-2}+e_t,\qquad H_0\!:\theta=0\ \text{vs}\ H_1\!:\theta<0 . \end{equation}
\noindent\textit{Code:}
\begingroup\footnotesize
\begin{verbatim}
d = wooldridge.data("traffic2")
def df_reg(y, lags, trend):                       # OLS of dy on const, y_{t-1}, dy lags (, t)
    y = pd.Series(y); dy = y.diff(); X = pd.DataFrame({"ylag": y.shift(1)})
    for j in range(1, lags+1): X[f"dl{j}"] = dy.shift(j)
    if trend: X["t"] = np.arange(1, len(y)+1)
    D = pd.concat([dy.rename("dy"), sm.add_constant(X)], axis=1).dropna()
    return sm.OLS(D.dy, D.drop(columns="dy")).fit()
\end{verbatim}
\endgroup
\noindent\textit{Key output (coefficient on $y_{t-1}$):}
\begin{center}\small
\begin{tabular}{llcccccc}
\toprule
Series & Regression & $n$ & $\hat\theta$ & s.e. & $t_{\hat\theta}$ & 2.5\% c.v. & Reject at 2.5\%?\\
\midrule
$ltotacc$ & (i) DF, no lags & 107 & $-0.1923$ & 0.0581 & $-3.31$ & $-3.12$ & yes\\
$ltotacc$ & (ii) ADF, 2 lags & 105 & $-0.0850$ & 0.0566 & $-1.50$ & $-3.12$ & no\\
$ltotacc$ & (iii) ADF, 2 lags + trend & 105 & $-0.4253$ & 0.1160 & $-3.67$ & $-3.66$ & yes (barely)\\
$prcfat$ & DF, no lags & 107 & $-0.2906$ & 0.0689 & $-4.22$ & $-3.12$ & yes\\
$prcfat$ & (v) ADF, 2 lags & 105 & $-0.3837$ & 0.0809 & $-4.74$ & $-3.12$ & yes\\
$prcfat$ & (v) ADF, 2 lags + trend & 105 & $-0.4581$ & 0.0866 & $-5.29$ & $-3.66$ & yes\\
\bottomrule
\end{tabular}
\end{center}

\passo{(i)}{From \eqref{q13:adf} with $\delta=\gamma_1=\gamma_2=0$: $t_{\hat\theta}=-0.1923/0.0581=-3.31<-3.12$. The unit root is rejected at 2.5\%.}
\passo{(ii)}{With two lagged changes, $t_{\hat\theta}=-0.0850/0.0566=-1.50>-2.57$: we cannot reject a unit root
even at 10\%. The lagged changes are highly significant ($t=-4.75$ and $-3.06$), so (i) omitted relevant dynamics
and its rejection is not reliable. The ADF result is the one to trust.}
\passo{(iii)}{Adding $t$: $t_{\hat\theta}=-0.4253/0.1160=-3.67<-3.66$, a (marginal) rejection at 2.5\% against the
trend critical value. The trend is significant: $\hat\delta=0.0013$ with $t=3.31$, about $0.13\%$ more accidents per month.}
\passo{(iv)}{Without a trend, the slow upward drift in accidents makes $ltotacc_t$ look like a unit-root
process. Once the linear trend is allowed, the unit root is rejected. The best characterization is an
\emph{I(0) process about a linear time trend} (trend-stationary), with the caveat that the rejection is borderline.}
\passo{(v)}{For $prcfat_t$ with two lags: $t_{\hat\theta}=-4.74$ without a trend ($<-3.43$, reject at 1\%) and
$-5.29$ with a trend ($<-3.96$, reject at 1\%). It does not matter whether the trend is included: $prcfat_t$
is I(0) in both cases.}
\intuicao{A Dickey--Fuller test asks whether shocks die out. Omitted short-run dynamics bias it, and so does
an omitted deterministic trend: a trending stationary series looks like a random walk with drift if the test
cannot ``see'' the trend. The fatality percentage is a ratio with no trend and mean reversion, so the
specification does not matter for it.}
\refbib{Wooldridge (2016) \S18-2 (pp.~574--578, Eqs.~(18.21), (18.24), Tables 18.2--18.3); Tsay (2010) \S2.7.3 and \S2.7.5 (pp.~75--80, trend-stationary series and unit-root tests).}
\end{sol}

% ---------------------------------------------------------------- Q14
\questao{Wooldridge, Problem 12.4 (ARCH and serial correlation).}
\fonte{Wooldridge (2016), \textit{Introductory Econometrics: A Modern Approach}, 6th ed., Chapter 12,
Problem 4, printed p.~397, reproduced verbatim.}

\noindent\textbf{4} True or false: ``If the errors in a regression model contain ARCH, they must be serially correlated.''
\resolespaco
\begin{sol}
\begin{notacao}
\item $u_t$: regression error; $F_{t-1}$: information at $t-1$ (past errors and regressors).
\item ARCH(1), Wooldridge Eq.~(12.49): $E(u_t^2\mid F_{t-1})=\alpha_0+\alpha_1u_{t-1}^2$, $\alpha_0>0$, $0\le\alpha_1<1$; and $E(u_t\mid F_{t-1})=0$.
\end{notacao}
\noindent\textbf{Answer: false.}
\passo{zero conditional mean}{ARCH restricts only the conditional \emph{variance}. The error keeps}
\begin{equation}\label{q14:m} E(u_t\mid F_{t-1})=0 . \end{equation}
\passo{no autocorrelation}{For $j\ge1$, $u_{t-j}$ belongs to $F_{t-1}$. By the law of iterated expectations,}
\begin{align*}
\operatorname{Cov}(u_t,u_{t-j}) &= E(u_tu_{t-j})\\
&= E\big[u_{t-j}\,E(u_t\mid F_{t-1})\big]\\
&= E\big[u_{t-j}\cdot(0)\big] \qquad\text{using \eqref{q14:m}}\\
&= 0 .
\end{align*}
\passo{but dependence}{The errors are not independent: their squares are autocorrelated. Writing
$u_t^2=\alpha_0+\alpha_1u_{t-1}^2+v_t$ with $E(v_t\mid F_{t-1})=0$ (Wooldridge Eq.~(12.50)), $u_t^2$ is an AR(1)
with $\operatorname{Corr}(u_t^2,u_{t-1}^2)=\alpha_1$. A simulated ARCH(1) with $\alpha_0=1$, $\alpha_1=0.5$ and
$n=200{,}000$ (check script) gives $\widehat{\operatorname{Corr}}(u_t,u_{t-1})=-0.003$ and
$\widehat{\operatorname{Corr}}(u_t^2,u_{t-1}^2)=0.458$. The population value is $\alpha_1=0.5$; it exists because the
fourth moment is finite ($3\alpha_1^2=0.75<1$). The sample value falls short because the eighth moment is infinite
($105\alpha_1^4=6.56\ge1$), so the sample autocorrelation of $u_t^2$ converges slowly.}
\intuicao{ARCH errors are a martingale difference sequence: unpredictable in level, predictable in size.
OLS standard errors that assume no serial correlation remain valid (heteroskedasticity-robust ones are
needed), and tests for serial correlation should not be confused with tests for ARCH (Ljung--Box on $u_t$ vs on $u_t^2$).}
\refbib{Wooldridge (2016) \S12-6c (pp.~393--394, Eqs.~(12.49)--(12.50)); Tsay (2010) \S3.4.1 (pp.~117--118, properties of ARCH models: $a_t$ uncorrelated, $a_t^2$ autocorrelated).}
\end{sol}

\end{document}
